% Half-title, title and copyright pages.
%
% \bookversion is the release the PDF was built for: `make VERSION=v1.2.3` passes it in, and the
% release workflow passes the tag. Without it the title page carries only the date of the build.
\providecommand{\bookversion}{}
\newcommand{\bookdate}{\number\day\space\ifcase\month\or January\or February\or March\or April\or
  May\or June\or July\or August\or September\or October\or November\or December\fi\space
  \number\year}
\newcommand{\bookedition}{\ifx\bookversion\empty\else\bookversion\,·\,\fi\bookdate}

\thispagestyle{empty}
\vspace*{2.2in}
\begin{flushright}
  {\Large\bfseries Zephyr RTOS}
\end{flushright}
\cleardoublepage

\thispagestyle{empty}
\vspace*{1.1in}
\begin{flushleft}
  {\color{accent}\rule{0.9in}{2.5pt}}\par\vspace{1.4ex}
  {\fontsize{32}{38}\selectfont\bfseries Zephyr\\[0.15ex] RTOS\par}
  \vspace{2.2ex}
  {\Large\itshape The build, the kernel and a driver,\\on a board that does not exist\par}
  \vspace{1.2in}
  {\large Erik Pihl\par}
  \vfill
  {\small\scshape\lsstyle qrtech academy\par}
  \vspace{0.6ex}
  {\small\color{muted} \bookedition\par}
\end{flushleft}
\newpage

\thispagestyle{empty}
\vspace*{\fill}
{\small
\noindent\textit{Zephyr RTOS}\\
Copyright \textcopyright\ 2026 QRTECH Academy.

\medskip
\noindent The text and figures of this book, and of the course it is typeset from, are licensed
under the Creative Commons Attribution-NonCommercial-ShareAlike 4.0 International License
(CC BY-NC-SA 4.0). You may copy and redistribute them, and adapt them, for any non-commercial
purpose, provided you give appropriate credit and share what you make under the same license. To
view a copy of the license, visit \url{https://creativecommons.org/licenses/by-nc-sa/4.0/}.

\medskip
\noindent The code in this book may also be used under the MIT License, which is the license of all
the code in the companion repository. The repository holds the lecture material, the example
programs, the exercises and their test suites:

\medskip
\noindent\url{https://github.com/qrtech-academy/zephyr}

\medskip
\noindent Every program in this book is built for Zephyr 4.4.2 and the board
\code{native_sim/native/64}, with the host's GCC, on Linux. Zephyr itself is licensed under
Apache~2.0 and is not part of this book; \secref{c1:app:a} fetches it. Zephyr, Linux and the other
product names in this book are trademarks of their respective owners.

\medskip
\noindent Set in TeX Gyre Pagella and DejaVu Sans Mono with Lua\LaTeX.

\medskip
\noindent\textcolor{muted}{Typeset from the six lectures of the course and their exercises. This
edition: \bookedition.}
\par}
\cleardoublepage
